
\paragraph{Example} We conclude this section with a last example showing how the tensor network graphical notation is very useful for representing complicated operations on tensors:
\begin{align}\label{eq:einsum}
\begin{tikzpicture}[baseline = \baseline]
    \tikzset{tensor/.style = {minimum size = 0.5cm,shape = circle,thick,draw=black,inner sep = 0pt}, edge/.style = 
    {thick,line width=.4mm},every loop/.style={}}
    \def\y{-0.5}
    \def\x{0}
    \node[tensor,fill=green!50!red!50!white] (At) at (\x,\y){\scalebox{0.85}{$\At$}};
    \node[tensor,fill=red!20!white] (S) at (\x-1,\y){\scalebox{0.85}{$\St$}};
    \node[tensor,fill=red!50!white] (Tt) at (\x,\y+1){\scalebox{0.85}{$\Tt$}};
    \node[tensor,fill=blue!40!red!60!white] (A) at (\x+1,\y){\scalebox{0.85}{$\Ab$}};
    \draw[edge] (Tt) -- (At) node [midway,left] {\scalebox{0.7}{\dimcolor{$R$}}};
    \draw[edge] (Tt) -- (A) node [above,midway] {\scalebox{0.7}{\dimcolor{$R$}}};
    \draw[edge] (At) -- (A)  node [midway,above] {\scalebox{0.7}{\dimcolor{$R$}}};
    \draw[edge] (Tt) -- (\x,\y+1.75) node [pos=1.3] {\scalebox{0.7}{\indcolor{$k$}}};
    \draw[edge] (At) -- (S)  node [midway,above] {\scalebox{0.7}{\dimcolor{$R$}}};
    \draw[edge] (S) -- (Tt)  node [midway,above] {\scalebox{0.7}{\dimcolor{$R$}}};
    \draw[edge] (Tt) -- (\x-0.75,\y+1)  node [pos=1.3] {\scalebox{0.7}{\indcolor{$j$}}};
    \draw[edge] (Tt) -- (\x+0.75,\y+1)  node [pos=1.3] {\scalebox{0.7}{\indcolor{$l$}}};
    \draw[edge] (S) -- (\x-1.2,\y+0.75)  node [pos=1.3] {\scalebox{0.7}{\indcolor{$i$}}};
    \end{tikzpicture}
&=
\sum_{r_1=1}^R\sum_{r_2=1}^R\sum_{r_3=1}^R\sum_{r_4=1}^R\sum_{r_5=1}^R\St_{ir_1 r_2}\At_{r_2r_3r_5}\Tt_{iklr_1r_3r_4}\Ab_{r_4r_5}\nonumber\\
&=
    \begin{tikzpicture}[baseline = \baseline]
    \tikzset{tensor/.style = {minimum size = 0.5cm,shape = circle,thick,draw=black,inner sep = 0pt}, edge/.style = 
    {thick,line width=.4mm},every loop/.style={}}
    \def\y{0}
    \def\x{0}
    \node[tensor,fill=blue!50!red!30!yellow!50!white] (Ht) at (\x,\y){\scalebox{0.85}{$\Ht$}};
    \draw[edge](Ht) -- (\x,\y+0.7) node [pos=1.3] {\scalebox{0.7}{\indcolor{$j$}}};
    \draw[edge](Ht) -- (\x,\y-0.7) node [pos=1.3] {\scalebox{0.7}{\indcolor{$l$}}};
    \draw[edge](Ht) -- (\x-0.7,\y) node [pos=1.3] {\scalebox{0.7}{\indcolor{$i$}}};
    \draw[edge](Ht) -- (\x+0.7,\y) node[pos=1.3] {\scalebox{0.7}{\indcolor{$k$}}};
    \end{tikzpicture}.
\end{align}
\subsection{Copy Tensors~(Hyperedges)}\label{sec:copy-tensor}


Sometimes, one needs to represent a simultaneous contraction shared among more than two tensors. Consider, for example, the following contraction between three vectors, $\ab\in\RR^d, \bb\in\RR^d,$ and $\cbb\in\RR^d$ to obtain a scalar: $\sum_{i=1}^d\ab_i\bb_i\cbb_i$.
This operation can be depicted in tensor networks using a special tensor called \emph{copy tensor}. The copy tensor is equivalent to a Kronecker delta, i.e., a hyper-diagonal tensor with ones on the diagonal and zeros elsewhere, and is represented as a black dot in tensor network diagrams. E.g., the copy tensor  
\begin{tikzpicture}[baseline=\baseline]
   \tikzset{tensor/.style = {minimum size = 0.5cm,shape = circle,thick,draw=black,inner sep = 0pt}, edge/.style = {thick,line width=.4mm},every loop/.style={}}
    \def\y{0}
    \def\x{0}
     \draw  node[fill = black,circle,inner sep=0pt,minimum size=4pt] (m) at (\x,\y) {};
     \draw  node[draw=none] (a) at (\x-0.7,\y) {};
     \draw  node[draw=none] (b) at (\x+0.7,\y) {};
     \draw  node[draw=none] (c) at (\x,\y-0.7) {};
    \draw[edge] (a) -- (m);
    \draw[edge] (m) -- (b);
    \draw[edge] (m) -- (c);
\end{tikzpicture}
is defined element-wise as
$$
 \left(\begin{tikzpicture}[baseline=\baseline]
   \tikzset{tensor/.style = {minimum size = 0.5cm,shape = circle,thick,draw=black,inner sep = 0pt}, edge/.style = {thick,line width=.4mm},every loop/.style={}}
    \def\y{0}
    \def\x{0}
     \draw  node[fill = black,circle,inner sep=0pt,minimum size=4pt] (m) at (\x,\y) {};
     \draw  node[draw=none] (a) at (\x-0.7,\y) {};
     \draw  node[draw=none] (b) at (\x+0.7,\y) {};
     \draw  node[draw=none] (c) at (\x,\y-0.7) {};
    \draw[edge] (a) -- (m);
    \draw[edge] (m) -- (b);
    \draw[edge] (m) -- (c);
\end{tikzpicture}\right)_{ijk} = \delta_{ijk} = \begin{cases}
      1 & \text{if } i=j=k\\
      0 & \text{otherwise}
    \end{cases}.
$$
Using the copy tensor, the 3-way contraction mentioned above can be represented as 
\begin{align*}
\begin{tikzpicture}[baseline=\baseline]
   \tikzset{tensor/.style = {minimum size = 0.5cm,shape = circle,thick,draw=black,inner sep = 0pt}, edge/.style = {thick,line width=.4mm},every loop/.style={}}
    \def\y{0}
    \def\x{0}
    \node[tensor,fill=blue!40!red!60!white] (a) at (\x,\y){\scalebox{0.85}{$\ab$}};
    \node[tensor,fill=blue!10!white] (b) at (\x+1.5,\y){\scalebox{0.85}{$\bb$}};
    \node[tensor,fill=blue!20!red!40!white] (c) at (\x+0.7,\y-0.75){\scalebox{0.85}{$\cbb$}};
    \draw  node[fill = black,circle,inner sep=0pt,minimum size=4pt] (m) at (\x+0.7,\y) {};
    \draw[edge] (a) -- (m) node [above,midway] {\scalebox{0.7}{\dimcolor{$d$}}};
    \draw[edge] (m) -- (b) node [above,midway] {\scalebox{0.7}{\dimcolor{$d$}}};
    \draw[edge] (m) -- (c) node [left,midway] {\scalebox{0.7}{\dimcolor{$d$}}};
    \end{tikzpicture} = \sum_{i,j,k=1}^d\delta_{ijk}\ab_i\bb_j\cbb_k =
\sum_{i=1}^d\ab_i\bb_i\cbb_i .
\end{align*}
We now give a more formal definition of the copy tensor of arbitrary order $N$.
\begin{definition}
N-th order copy tensor
   \begin{tikzpicture}[baseline=\baseline]
    \tikzset{tensor/.style = {minimum size = 0.5cm,shape = circle,thick,draw=black,inner sep = 0pt}, edge/.style   = {thick,line width=.4mm},every loop/.style={}}
    \def\y{0}
    \def\x{0}
    \draw  node[fill = black,circle,inner sep=0pt,minimum size=4pt] (m) at (\x,\y) {};
    \draw[edge] (m) -- (\x-0.6,0) node [midway,above] {\scalebox{0.7}{\dimcolor{$d$}}};
    \draw[edge] (m) -- (\x+0.6,0) node [midway,above] {\scalebox{0.7}{\dimcolor{$d$}}};;
    \draw[edge] (m) -- (\x-0.5,-0.3) node [midway, below] {\scalebox{0.7}{\dimcolor{$d$}}};;
    \draw[dotted] (m) -- (\x,-0.6);
    \draw[dotted] (m) -- (\x,-0.6);
    \draw[dotted] (m) -- (\x+0.5,-0.5);
    \draw[dotted] (m) -- (\x+0.3,-0.6);
\end{tikzpicture}
is the tensor of shape $\underbrace{d\times d\dots\times d}_{N\text{ times}}$ defined by 
\begin{align*}
       \begin{tikzpicture}[baseline=\baseline]
    \tikzset{tensor/.style = {minimum size = 0.5cm,shape = circle,thick,draw=black,inner sep = 0pt}, edge/.style   = {thick,line width=.4mm},every loop/.style={}}
    \def\y{0}
    \def\x{0}
    \draw  node[fill = black,circle,inner sep=0pt,minimum size=4pt] (m) at (\x,\y) {};
    \draw[edge] (m) -- (\x-0.6,0) node [midway,above] {\scalebox{0.7}{\dimcolor{$d$}}};
    \draw[edge] (m) -- (\x+0.6,0) node [midway,above] {\scalebox{0.7}{\dimcolor{$d$}}};;
    \draw[edge] (m) -- (\x-0.5,-0.3) node [midway, below] {\scalebox{0.7}{\dimcolor{$d$}}};;
    \draw[dotted] (m) -- (\x,-0.6);
    \draw[dotted] (m) -- (\x,-0.6);
    \draw[dotted] (m) -- (\x+0.5,-0.5);
    \draw[dotted] (m) -- (\x+0.3,-0.6);
\end{tikzpicture}
= \sum_{i=1}^d\underbrace{\eb_i\outprod\eb_i\outprod\dots\outprod\eb_i}_{N\text{ times}},%
\end{align*}
where $\eb_1,\eb_2,\dots,\eb_d$ are the vectors of the canonical basis of $\RR^{d}$.
\end{definition}
The name copy tensor comes from the fact that one can "copy" the canonical vectors by contraction with the copy tensor. Indeed, let $\eb_i\in\RR^{d}$ be a canonical basis vector, then one 
can check that
$       
\begin{tikzpicture}[baseline=\baseline]
    \tikzset{tensor/.style = {minimum size = 0.5cm,shape = circle,thick,draw=black,inner sep = 0pt}, edge/.style   = {thick,line width=.4mm},every loop/.style={}}
    \def\y{0}
    \def\x{0}
    \node[tensor,fill=blue!40!red!60!white] (e) at (\x,\y){\scalebox{0.85}{$\eb_i$}};
    \draw  node[fill = black,circle,inner sep=0pt,minimum size=4pt] (m) at (\x+0.75,\y) {};
    \draw[edge] (e) -- (m) node [midway,above] {\scalebox{0.7}{\dimcolor{$d$}}};
    \draw[edge] (m) -- (\x+1.4,\y) node [midway,above] {\scalebox{0.7}{\dimcolor{$d$}}};
    \draw[edge] (m) -- (\x+0.75,\y-0.7) node [midway,right] {\scalebox{0.7}{\dimcolor{$d$}}};
    \node[draw=none] () at (\x+2,\y) {\scalebox{1}{$=$}};
    \node[tensor,fill=blue!40!red!60!white] (e1) at (\x+2.5,\y){\scalebox{0.85}{$\eb_i$}};
    \node[tensor,fill=blue!40!red!60!white] (e2) at (\x+3.5,\y){\scalebox{0.85}{$\eb_i$}};
    \draw[edge] (e2) -- (\x+3.5,\y-0.7) node [midway,right] {\scalebox{0.7}{\textcolor{gray}{$d$}}};
    \draw[edge] (e1) -- (\x+2.5,\y-0.7) node [midway,right] {\scalebox{0.7}{\textcolor{gray}{$d$}}};
\end{tikzpicture}.
$
Here is a short proof of this fact:
\begin{align*}
    \left(
\begin{tikzpicture}[baseline=\baseline]
    \tikzset{tensor/.style = {minimum size = 0.5cm,shape = circle,thick,draw=black,inner sep = 0pt}, edge/.style   = {thick,line width=.4mm},every loop/.style={}}
    \def\y{0}
    \def\x{0}
    \node[tensor,fill=blue!40!red!60!white] (e) at (\x,\y){\scalebox{0.85}{$\eb_i$}};
    \draw  node[fill = black,circle,inner sep=0pt,minimum size=4pt] (m) at (\x+0.75,\y) {};
    \draw[edge] (e) -- (m) node [midway,above] {\scalebox{0.7}{\dimcolor{$d$}}};
    \draw[edge] (m) -- (\x+1.4,\y) node [midway,above] {\scalebox{0.7}{\dimcolor{$d$}}};
    \draw[edge] (m) -- (\x+0.75,\y-0.7)node [midway,right] {\scalebox{0.7}{\dimcolor{$d$}}};
\end{tikzpicture}\right)_{jk}
&= 
\sum_{s}(\eb_i)_{s}\delta_{sjk} = \sum_{s}\delta_{is}\delta_{sjk} = \delta_{ijk}\\ 
&= 
\begin{cases}
     1 & \text{if } i=j=k\\
      0 & \text{otherwise}.
    \end{cases} =
    (\eb_i\eb_i^\ts)_{j,k} =
\begin{tikzpicture}[baseline=\baseline]
    \tikzset{tensor/.style = {minimum size = 0.5cm,shape = circle,thick,draw=black,inner sep = 0pt}, edge/.style   = {thick,line width=.4mm},every loop/.style={}}
    \def\y{0}
    \def\x{0}
    \node[tensor,fill=blue!40!red!60!white] (e1) at (\x,\y){\scalebox{0.85}{$\eb_i$}};
    \node[tensor,fill=blue!40!red!60!white] (e2) at (\x+1,\y){\scalebox{0.85}{$\eb_i$}};
    \draw[edge] (e2) -- (\x+1,\y-0.7) node [below] {\scalebox{0.7}{\indcolor{$k$}}};
    \draw[edge] (e1) -- (\x,\y-0.7) node [below] {\scalebox{0.7}{\indcolor{$j$}}};
\end{tikzpicture}.
\end{align*}
It is important to observe that the "copy" property of the copy tensor only holds for a canonical basis vector; it is not true for an arbitrary vector: in general
$       
\begin{tikzpicture}[baseline=\baseline]
    \tikzset{tensor/.style = {minimum size = 0.5cm,shape = circle,thick,draw=black,inner sep = 0pt}, edge/.style   = {thick,line width=.4mm},every loop/.style={}}
    \def\y{0}
    \def\x{0}
    \node[tensor,fill=blue!40!red!60!white] (e) at (\x,\y){\scalebox{0.85}{$\vb$}};
    \draw  node[fill = black,circle,inner sep=0pt,minimum size=4pt] (m) at (\x+0.75,\y) {};
    \draw[edge] (e) -- (m) node [midway,above] {\scalebox{0.7}{\dimcolor{$d$}}};
    \draw[edge] (m) -- (\x+1.4,\y) node [midway,above] {\scalebox{0.7}{\dimcolor{$d$}}};
    \draw[edge] (m) -- (\x+0.75,\y-0.7) node [midway,right] {\scalebox{0.7}{\dimcolor{$d$}}};
    \node[draw=none] () at (\x+2,\y) {\scalebox{1}{$\neq$}};
    \node[tensor,fill=blue!40!red!60!white] (e1) at (\x+2.5,\y){\scalebox{0.85}{$\vb$}};
    \node[tensor,fill=blue!40!red!60!white] (e2) at (\x+3.5,\y){\scalebox{0.85}{$\vb$}};
    \draw[edge] (e2) -- (\x+3.5,\y-0.7) node [midway,right] {\scalebox{0.7}{\dimcolor{$d$}}};
    \draw[edge] (e1) -- (\x+2.5,\y-0.7) node [midway,right] {\scalebox{0.7}{\dimcolor{$d$}}};
\end{tikzpicture}
$. Since the copy tensor only "copies" basis vectors because $\delta_{ijk}$ picks out diagonal entries.
\begin{proposition}\label{prop:copy.tensor} In the following, we list some useful properties of copy tensors.
\begin{enumerate}
    \item The simplest version of the copy tensor is an all-ones vector, i.e., 
$$
  \begin{tikzpicture}[baseline=\baseline]
    \tikzset{tensor/.style = {minimum size = 0.5cm,shape = circle,thick,draw=black,inner sep = 0pt}, edge/.style   = {thick,line width=.4mm},every loop/.style={}}
    \def\y{0.5}
    \def\x{0}
    \draw  node[fill = black,circle,inner sep=0pt,minimum size=4pt] (m) at (\x,\y) {};
    \draw[edge] (m) -- (\x,\y-1) node [left,midway] {\scalebox{0.7}{\dimcolor{$n$}}};
 \end{tikzpicture}
 =\sum_{i=1}^n\eb_i =\overrightarrow{\mathbf{1}}.
$$\label{itm:one-copy-tensor}
\item The second-order copy tensor is the identity matrix, where we can simply omit the black node in the middle and represent it simply as a line, i.e.,  
$
    \begin{tikzpicture}[baseline=\baseline]
    \tikzset{tensor/.style = {minimum size = 0.5cm,shape = circle,thick,draw=black,inner sep = 0pt}, edge/.style   = {thick,line width=.4mm},every loop/.style={}}
    \def\y{0}
    \def\x{0}
    \draw  node[fill = black,circle,inner sep=0pt,minimum size=4pt] (m) at (\x,\y) {};
    \draw  node[fill = black,circle,inner sep=0pt,minimum size=4pt] (m) at (\x,\y) {};
    \draw[edge] (m) -- (\x-0.6,0) node [very near end,left] {\scalebox{0.7}{\indcolor{$i$}}};
    \draw[edge] (m) -- (\x+0.6,0) node [very near end,right]
    {\scalebox{0.7}{\indcolor{$i$}}};
    \node[draw=none] at (\x+1,\y) {\scalebox{1}{\textcolor{black}{$=$}}};;
    \draw[edge] (\x+1.3,\y) -- (\x+2.5,\y) node [very near start,left=0.05cm] {\scalebox{0.7}{\indcolor{$i$}}} node [very near end,right=0.05cm] {\scalebox{0.7}{\indcolor{$i$}}};
    \end{tikzpicture} = \delta_{ij} = \begin{cases}
      1 & \text{if } i=j\\
      0 & \text{otherwise},
    \end{cases}
    =
    \Ib_{ij}.
$
\item 
Interestingly, when contracting any number of legs from one copy tensor to legs of another copy tensor, the resulting tensor network is itself a copy
tensor (obtained by merging the contracted legs and dots in the tensor network diagram). For example: 
\begin{align*}  
\sum_{l_1}\delta_{i_1j_1k_1l_1}\delta_{i_2j_2k_2l_1} 
&=
\begin{tikzpicture}[baseline=\baseline]
    \tikzset{tensor/.style = {minimum size = 0.5cm,shape = circle,thick,draw=black,inner sep = 0pt}, edge/.style   = {thick,line width=.4mm},every loop/.style={}}
    \def\y{0}
    \def\x{0}
    \draw  node[fill = black,circle,inner sep=0pt,minimum size=4pt] (m) at (\x,\y) {};
    \draw[edge] (m) -- (\x,\y-0.5)node [below] {\scalebox{0.6}{\indcolor{$i_1$}}};
    \draw[edge](m) -- (\x,\y+0.5)node [above] {\scalebox{0.6}{\indcolor{$
    k_1$}}};
    \draw[edge] (m) -- (\x-0.5,\y)node [left] {\scalebox{0.6}{\indcolor{$j_1$}}};
    \draw  node[fill = black,circle,inner sep=0pt,minimum size=4pt] (n) at (\x+1,\y) {};
    \draw[edge] (n) -- (\x+1.5,\y)node [right] {\scalebox{0.6}{\indcolor{$j_2$}}};
    \draw[edge] (n) --  (\x+1,\y+0.5)node [above] {\scalebox{0.6}{\indcolor{$k_2$}}};   
    \draw[edge] (n) -- (\x+1,\y-0.5)node [below] {\scalebox{0.6}{\indcolor{$i_2$}}};
    \draw[edge](m) -- (n)node [midway, above] {\scalebox{0.6}{\indcolor{$l_1$}}};
 \end{tikzpicture}
= 
\delta_{i_1j_1k_1i_2j_2k_2}\\
&=
\begin{tikzpicture}[baseline=\baseline]
    \tikzset{tensor/.style = {minimum size = 0.5cm,shape = circle,thick,draw=black,inner sep = 0pt}, edge/.style   = {thick,line width=.4mm},every loop/.style={}}
    \def\y{0}
    \def\x{0}
    \draw node[fill = black,circle,inner sep=0pt,minimum size=5pt] (m) at (\x,\y) {};
    \draw[edge] (m) -- (\x+0.75,\y)node [right] {\scalebox{0.6}{\indcolor{$j_2$}}};
    \draw[edge] (m) -- (\x-0.75,\y)node [left] {\scalebox{0.6}{\indcolor{$j_1$}}};
    \draw[edge] (m) -- (\x+0.5,\y-0.5)node [below] {\scalebox{0.6}{\indcolor{$i_2$}}};
    \draw[edge] (m) -- (\x-0.5,\y+0.5)node [above] {\scalebox{0.6}{\indcolor{$k_1$}}};
    \draw[edge] (m) -- (\x-0.5,\y-0.5)node [below] {\scalebox{0.6}{\indcolor{$i_1$}}};
    \draw[edge] (m) -- (\x+0.5,\y+0.5)node [above] {\scalebox{0.6}{\indcolor{$k_2$}}};
\end{tikzpicture}
=
 \begin{tikzpicture}[baseline=\baseline]
    \tikzset{tensor/.style = {minimum size = 0.5cm,shape = circle,thick,draw=black,inner sep = 0pt}, edge/.style   = {thick,line width=.4mm},every loop/.style={}}
    \def\y{0}
    \def\x{0}
    \draw  node[fill = black,circle,inner sep=0pt,minimum size=5pt] (m) at (\x,\y+0.4) {};
    \draw  node[fill = black,circle,inner sep=0pt,minimum size=5pt] (n) at (\x,\y-0.4) {};
    \draw[edge] (n) -- (m)node [midway, right] {\scalebox{0.6}{\indcolor{$i_1$}}};;
    \draw[edge](m) -- (\x,\y+1)node [above] {\scalebox{0.6}{\indcolor{$k_2$}}};
    \draw[edge](n) -- (\x,\y-1)node [below] {\scalebox{0.6}{\indcolor{$k_1$}}};
    \draw[edge](m) -- (\x+0.45,\y+0.65)node [above] {\scalebox{0.6}{\indcolor{$l_2$}}};
    \draw[edge](m) -- (\x-0.45,\y+0.65)node [above] {\scalebox{0.6}{\indcolor{$j_2$}}};
    \draw[edge](n) -- (\x+0.45,\y-0.65)node [above] {\scalebox{0.6}{\indcolor{$l_1$}}};;
    \draw[edge](n) -- (\x-0.45,\y-0.65)node [above] {\scalebox{0.6}{\indcolor{$j_1$}}};;
    \end{tikzpicture}
=
\sum_{i_1}\delta_{i_1j_1k_1l_1}\delta_{i_1j_2k_2l_2}.
\end{align*}
\item\label{prop:diagcopy}
For any vector $\vb\in\RR^n$, let $\diag(\vb)\in\RR^{n\times n}$ denote the diagonal matrix having the entries of $\vb$ on the diagonal, and for any matrix $\Ab\in\RR^{n\times n}$, let $\diag(\Ab)\in\RR^{n}$ denote the vector containing the diagonal entries of $\Ab\in\RR^{n\times n}$.
